\subsection{Reality}
\label{pg:sec:reality}

The action must be real.
For the $\Psi$ sector this is standard: with the closed-string reality condition $\Psi^\ddagger=\Psi$, the kinetic term $\frac12\braket{\Psi|c_0^-Q_B|\Psi}$ is real \cite{Erbin:2021:StringFieldTheory}.
This section does the same for the $\Sigma$ sector of the quartic action~\eqref{p:eq:hs4:S4tilde}: the conjugation of the pairing, of the vertex and of $p_{2,2}$, the condition on $\Sigma$ that follows, and the sense in which that condition is forced rather than chosen.
The argument is algebraic and rests on one hypothesis about the vertex, stated where it is used.

\subsubsection{The conjugation}
\label{pg:sec:reality:conj}

The conjugation $\ddagger$ is the composition of BPZ and Hermitian conjugation.
It is an antilinear involution of $\mc H_0$ preserving the ghost number: $(aA+bB)^\ddagger=a^*A^\ddagger+b^*B^\ddagger$ and $(A^\ddagger)^\ddagger=A$.
A state is \emph{real} if $A^\ddagger=A$; every state splits as $A=A_r+\I A_i$ with $A_r=\frac12(A+A^\ddagger)$ and $A_i=\frac1{2\I}(A-A^\ddagger)$ both real, so the real states span $\mc H_0$ over $\C$.
On two-string states we extend it componentwise,
\begin{equation}
  (A_1 \otimes A_2)^\ddagger
    := A_1^\ddagger \otimes A_2^\ddagger,
\end{equation}
with no reordering sign.
For $\Sigma$ this is unambiguous: both strands sit at $N_{\rm gh}=2$, even, so the alternative convention with reversed order, which differs by $(-1)^{k(k-1)/2}$ on tensors all of whose strands are odd, agrees with it.
$\Sym^2\mc H_0$ is $\ddagger$-stable, and $(\Psi^2)^\ddagger=\Psi^\ddagger\otimes\Psi^\ddagger=\Psi^2$.

A multilinear form $F$ is \emph{real} if it takes real values on real arguments.
This is equivalent to \emph{equivariance},
\begin{equation}
  \label{pg:eq:reality:equiv}
  F(A_1, \ldots, A_n)^*
    = F(A_1^\ddagger, \ldots, A_n^\ddagger)
  \quad\text{for all } A_i ,
\end{equation}
since both sides of \eqref{pg:eq:reality:equiv} are antilinear in each slot and agree on real arguments, which span.
We use the two formulations interchangeably.

\subsubsection{Conjugation of the ingredients}
\label{pg:sec:reality:ingredients}

\emph{The pairing.}
In the normalization $\braket{c_{-1}\bar c_{-1}c_0^-c_0^+c_1\bar c_1}=-\frac12$, real and without a factor of $\I$, the pairing $\braket{A,B}=\braket{A|c_0^-|B}$ is real on real states.
The only structural input is that Hermitian conjugation of $c_0^-$ and of the BPZ product introduces no phase, which is what the normalization fixes; it is the same choice that makes the $\Psi$ kinetic term real, and the reality of the $\Psi$ kinetic term already pairs $N_{\rm gh}=2$ with $N_{\rm gh}=3$.
One caveat: Hermitian conjugation reverses the order of $A$ and $B$, so for two odd arguments the equivariance of $\braket{\cdot,\cdot}$ may carry a relative sign $(-1)^{\epsilon(A)\epsilon(B)}$ that the even--odd case does not probe.
Every pairing below has one even argument, so the odd--odd case is never needed.
The two-string pairing is then real too: on product states $\braket{X,Y}_2$ is a product of two real pairings, both sides of the equivariance identity are antilinear in each argument, and product states span.

\emph{The vertex.}
The hypothesis the whole section rests on is that the quartic vertex is equivariant,
\begin{equation}
  \label{pg:eq:reality:V4}
  \mc V_4(A_1, A_2, A_3, A_4)^*
    = \mc V_4(A_1^\ddagger, A_2^\ddagger, A_3^\ddagger, A_4^\ddagger).
\end{equation}
This is the standard reality of a closed-string vertex, the one that makes the non-polynomial action real: it follows from the Hermiticity of the vertex operators and plumbing data, $\mc O(t^*)=\mc O(t)^*$ in the language of Sonoda and Zwiebach \cite{Sonoda:1990:CovariantClosedString}, together with the invariance of the moduli region of the vertex and of its measure under complex conjugation of the moduli.
We take \eqref{pg:eq:reality:V4} from the literature and do not re-derive it; everything below is conditional on it.

\emph{The product.}
Since $\braket{\cdot,\cdot}_2$ is real and non-degenerate, \eqref{pg:eq:reality:V4} says that $p_{2,2}$ is a real operator, $p_{2,2}(Y^\ddagger)=(p_{2,2}Y)^\ddagger$.
Indeed, conjugating \eqref{p:eq:conv:sprod} in two ways,
\begin{equation}
  \begin{aligned}
    \braket{X^\ddagger, (p_{2,2} Y)^\ddagger}_2
      &
      = \braket{X, p_{2,2} Y}_2^*
      = \mc V_4(X, Y)^*
    \\ &
      = \mc V_4(X^\ddagger, Y^\ddagger)
      = \braket{X^\ddagger, p_{2,2}(Y^\ddagger)}_2
    \quad\text{for all } X,
  \end{aligned}
\end{equation}
and non-degeneracy gives the claim.
A real operator has a $\ddagger$-stable kernel and image, and its inverse, where it exists, is real; in particular $\ker\mc V_4=\ker p_{2,2}$ is $\ddagger$-stable, so the quotient~\eqref{p:eq:hs4:quotient} inherits the conjugation and a reality condition on the class of $\Sigma$ makes sense.

\subsubsection{The reality condition on \texorpdfstring{$\Sigma$}{Sigma}}
\label{pg:sec:reality:condition}

Impose
\begin{equation}
  \label{pg:eq:reality:Sigma}
  \boxed{\;
  \Sigma^\ddagger = \Sigma .
  \;}
\end{equation}
Then both $\Sigma$ terms of \eqref{p:eq:hs4:S4tilde} are real: by \eqref{pg:eq:reality:V4}, $\mc V_4(\Sigma,\Sigma)^*=\mc V_4(\Sigma^\ddagger,\Sigma^\ddagger)=\mc V_4(\Sigma,\Sigma)$, and $\mc V_4(\Sigma,\Psi^2)^*=\mc V_4(\Sigma,\Psi^2)$ likewise, the coefficients being real.
The on-shell value $\Sigma=\Psi^2$ is real, and the class of a real solution in the quotient is real.
Reality is thus conditional on \eqref{pg:eq:reality:V4} and on nothing else; the two $c_0^-$ of $\braket{\cdot,\cdot}_2$ play no role beyond the reality of the pairing, and if a convention had attached a phase $\lambda$ to the single $c_0^-$ of $\braket{\cdot,\cdot}$, it would enter the $\Sigma$ kinetic term as $\lambda^2$.

Let's now ask whether \eqref{pg:eq:reality:Sigma} is forced.
The kinetic term alone does not force it: for $\Sigma=\I\Sigma_r$ with $\Sigma_r$ real, $\mc V_4(\Sigma,\Sigma)=-\mc V_4(\Sigma_r,\Sigma_r)$ is real as well, so $\Sigma^\ddagger=-\Sigma$ would also do.
The source term removes this freedom.
It is linear in $\Sigma$, and for $\Sigma=\Sigma_r+\I\Sigma_i$,
\begin{equation}
  \mc V_4(\Sigma, \Psi^2)
    = \mc V_4(\Sigma_r, \Psi^2)
      + \I \, \mc V_4(\Sigma_i, \Psi^2),
\end{equation}
which is real for all real $\Psi$ if and only if $\mc V_4(\Sigma_i,\Psi^2)=0$ for all real $\Psi$.
For a generic symmetric vertex the covectors $\mc V_4(\,\cdot\,,\Psi^2)$ span the dual of $\Sym^2\mc H_0$ and $\Sigma_i=0$; for the actual vertex, the directions annihilated by every $\mc V_4(\,\cdot\,,\Psi^2)$ are directions absent from the dynamics, and \eqref{pg:eq:reality:Sigma} is forced up to them.
A relative phase $\Sigma^\ddagger=\e^{\I\theta}\Sigma$ is excluded in the same way, and $\Sigma\to-\Sigma$ is a field redefinition that flips the sign of the source term, not a different reality condition.
So \eqref{pg:eq:reality:Sigma} is the unique condition compatible with a real $\Psi$, which is how the quartic step reads it off the equation of motion.

The argument was checked on a finite-dimensional model with a non-trivial antilinear involution, a random real reflection composed with complex conjugation, so that the real subspace is not the coordinate one, including a degenerate kernel and the two failure modes (a non-real component in a vertex, a relative phase on $\Sigma$).

